% ----------------------------------------------------
% Capítulo 1 - Introdução
% Edite o texto deste capítulo livremente.
% Use \undfsecao, \undfsubsecao e demais divisões conforme necessário.
% Este arquivo é incluído automaticamente no documento principal.
% ----------------------------------------------------

\chapter{Introdução}

O Trabalho de Conclusão de Curso constitui etapa relevante da formação acadêmica, pois reúne pesquisa, escrita científica e apresentação formal do conhecimento produzido. Neste modelo, o próprio TCC é tomado como tema de exemplo, com o objetivo de mostrar ao estudante uma estrutura inicial organizada e funcional.

A introdução deve apresentar o assunto do trabalho, delimitar o problema de pesquisa, indicar os objetivos e justificar a relevância do estudo. Também é recomendável antecipar, de forma breve, a organização geral dos capítulos.

\section{Contextualização do tema}

Escrever um TCC exige planejamento, leitura, revisão e capacidade de organizar ideias de maneira clara. Por isso, o domínio da estrutura do trabalho é parte importante da própria pesquisa.

\section{Objetivos}

Este modelo tem como objetivo servir de base para a redação de um TCC, oferecendo ao estudante um ponto de partida tecnicamente organizado e fácil de adaptar.